% =========================================================
% One-Sided Limits — Notes
% Chapter: functions
% Section: one-sided-limits
% Volume III · Analysis
% =========================================================

\section{One-Sided Limits}
\label{sec:one-sided-limits}

\begin{tcolorbox}[title={Toolkit --- One-Sided Limits}]
\begin{itemize}
  \item $\operatorname{RightLimit}(f,c,L)$: approach from $x>c$ only.
    Input condition: $0<x-c<\delta$.
  \item $\operatorname{LeftLimit}(f,c,L)$: approach from $x<c$ only.
    Input condition: $0<c-x<\delta$.
  \item Sequential criteria hold for each direction separately.
  \item $\operatorname{FunctionLimit}(f,c,L)\iff
    \operatorname{RightLimit}(f,c,L)\wedge\operatorname{LeftLimit}(f,c,L)$.
\end{itemize}
\end{tcolorbox}

% ---------------------------------------------------------
% Definition — One-Sided Limits
% ---------------------------------------------------------

\begin{tcolorbox}[title={Definition --- Right-Hand and Left-Hand Limits}]
\begin{definition}[Right-Hand and Left-Hand Limits]
\label{def:one-sided-limits}
Let $A\subseteq\mathbb{R}$, $f:A\to\mathbb{R}$, $c\in\mathbb{R}$.

(i) If $c$ is a cluster point of $A\cap(c,\infty)$,
$\operatorname{RightLimit}(f,c,L)$ holds if
\[
  \ForallOT{\varepsilon},\;
  \ExistsIT{\delta},\;
  \forall x\in A:\;
  0<x-c<\delta \Rightarrow \operatorname{OutputTolerance}(f(x),L,\varepsilon).
\]

(ii) If $c$ is a cluster point of $A\cap(-\infty,c)$,
$\operatorname{LeftLimit}(f,c,L)$ holds if
\[
  \ForallOT{\varepsilon},\;
  \ExistsIT{\delta},\;
  \forall x\in A:\;
  0<c-x<\delta \Rightarrow \operatorname{OutputTolerance}(f(x),L,\varepsilon).
\]
\end{definition}
\end{tcolorbox}

\begin{remark*}[Standard quantified statement]
\[
  \operatorname{RightLimit}(f,c,L)
  \;\iff\;
  \forall\varepsilon>0,\;\exists\delta>0,\;\forall x\in A:\;
  c<x<c+\delta \Rightarrow |f(x)-L|<\varepsilon.
\]
\end{remark*}

\begin{remark*}[Definition predicate reading]
\[
  \operatorname{RightLimit}(f,c,L)
  \;\colonequiv\;
  \ForallOT{\varepsilon},\;\ExistsIT{\delta},\;
  \forall x\in A\cap(c,c+\delta):\;
  \operatorname{OutputTolerance}(f(x),L,\varepsilon).
\]
\end{remark*}

\begin{remark*}[Negated quantified statement]
\[
  \neg\,\operatorname{RightLimit}(f,c,L)
  \;\iff\;
  \exists\varepsilon_0>0,\;\forall\delta>0,\;\exists x\in A:\;
  0<x-c<\delta \;\wedge\; |f(x)-L|\geq\varepsilon_0.
\]
\end{remark*}

\begin{remark*}[Negation predicate reading]
\[
  \neg\,\operatorname{RightLimit}(f,c,L)
  \;\colonequiv\;
  \exists\varepsilon_0>0,\;\forall\delta>0,\;\exists x\in A:\;
  0<x-c<\delta
  \;\wedge\;
  \neg\,\operatorname{OutputTolerance}(f(x),L,\varepsilon_0).
\]
\end{remark*}

\begin{remark*}[Interpretation]
$\operatorname{RightLimit}(f,c,L)$ is $\operatorname{FunctionLimit}$
restricted to $A\cap(c,\infty)$. The input condition $0<x-c<\delta$
is a one-sided instance of $\operatorname{InputTolerance}$. The left
limit is the symmetric restriction to $A\cap(-\infty,c)$.
\end{remark*}

% ---------------------------------------------------------
% Theorem — Sequential Criterion for Right-Hand Limits
% ---------------------------------------------------------

\begin{theorem}[Sequential Criterion for Right-Hand Limits]
\label{thm:sequential-criterion-right-hand-limit}
$\operatorname{RightLimit}(f,c,L)\iff$ for every
$(x_n)\subseteq A\cap(c,\infty)$ with $x_n\to c$,
$f(x_n)\to L$.

\smallskip\noindent
\hyperref[prf:sequential-criterion-right-hand-limit]{\textit{Go to proof.}}
\end{theorem}

\begin{remark*}[Standard quantified statement]
\[
  \operatorname{RightLimit}(f,c,L)
  \;\iff\;
  \forall(x_n)\subseteq A\cap(c,\infty):\;
  x_n\to c \Rightarrow
  \operatorname{ConvergentSequence}(f(x_n),L).
\]
\end{remark*}

\begin{remark*}[Definition predicate reading]
\[
  \operatorname{RightLimit}(f,c,L)\text{ (seq)}
  \;\colonequiv\;
  \forall(x_n)\subseteq A\cap(c,\infty):\;
  \operatorname{ConvergentSequence}(x_n,c)
  \Rightarrow
  \operatorname{ConvergentSequence}(f(x_n),L).
\]
\end{remark*}

% ---------------------------------------------------------
% Corollary — Sequential Criterion for Left-Hand Limits
% ---------------------------------------------------------

\begin{corollary}[Sequential Criterion for Left-Hand Limits]
\label{thm:sequential-criterion-left-hand-limit}
$\operatorname{LeftLimit}(f,c,L)\iff$ for every
$(x_n)\subseteq A\cap(-\infty,c)$ with $x_n\to c$,
$f(x_n)\to L$.

\smallskip\noindent
\hyperref[prf:sequential-criterion-left-hand-limit]{\textit{Go to proof.}}
\end{corollary}

\begin{remark*}[Standard quantified statement]
\[
  \operatorname{LeftLimit}(f,c,L)
  \;\iff\;
  \forall(x_n)\subseteq A\cap(-\infty,c):\;
  x_n\to c \Rightarrow
  \operatorname{ConvergentSequence}(f(x_n),L).
\]
\end{remark*}

\begin{remark*}[Symmetry]
The proof replaces $0<x-c<\delta$ with $0<c-x<\delta$ and
constructs $x_n$ with $c-1/n<x_n<c$ instead of $c<x_n<c+1/n$.
All other steps are identical.
\end{remark*}

% ---------------------------------------------------------
% Theorem — Two-Sided Limit via One-Sided Limits
% ---------------------------------------------------------

\begin{theorem}[Two-Sided Limit via One-Sided Limits]
\label{thm:two-sided-limit-iff-matching-one-sided-limits}
If $c$ is a cluster point of both $A\cap(c,\infty)$ and
$A\cap(-\infty,c)$, then
\[
  \operatorname{FunctionLimit}(f,c,L)
  \;\iff\;
  \operatorname{RightLimit}(f,c,L) \;\wedge\; \operatorname{LeftLimit}(f,c,L).
\]

\smallskip\noindent
\hyperref[prf:two-sided-limit-iff-matching-one-sided-limits]{\textit{Go to proof.}}
\end{theorem}

\begin{remark*}[Standard quantified statement]
\[
  \operatorname{FunctionLimit}(f,c,L)
  \;\iff\;
  \operatorname{RightLimit}(f,c,L) \;\wedge\; \operatorname{LeftLimit}(f,c,L).
\]
\end{remark*}

\begin{remark*}[Definition predicate reading]
\[
  \operatorname{FunctionLimit}(f,c,L)
  \;\colonequiv\;
  \operatorname{RightLimit}(f,c,L)
  \;\wedge\;
  \operatorname{LeftLimit}(f,c,L).
\]
\end{remark*}

\begin{remark*}[Negated quantified statement]
\[
  \neg\,\operatorname{FunctionLimit}(f,c,L)
  \;\iff\;
  \neg\,\operatorname{RightLimit}(f,c,L)
  \;\vee\;
  \neg\,\operatorname{LeftLimit}(f,c,L).
\]
\end{remark*}

\begin{remark*}[Failure mode]
$\operatorname{RightLimit}(f,c,L_+)\neq\operatorname{LeftLimit}(f,c,L_-)$
means the two-sided limit does not exist. Heaviside at $0$:
$H(0^-)=0\neq 1=H(0^+)$, so $\lim_{x\to 0}H(x)$ does not exist.
\end{remark*}

\begin{remark*}[Interpretation]
The two-sided limit exists iff both directions agree. The one-sided
limits are finer instruments detecting directional behavior; the
two-sided limit is a coarser summary requiring both to coincide.
\end{remark*}
